\documentclass[../main.tex]{subfiles}
\begin{document}
%==========================================TIÊU ĐỀ TẬP========================================
\begin{table}[h]
    \begin{adjustwidth}{-1cm}{}
        \begin{tabular}{>{\centering\arraybackslash}p{6cm}|>{\raggedright\arraybackslash}p{12.5cm}}
            \multirow{5}{6cm}
            {\\[-40pt]
                \flaregothic\fontsize{20pt}{20pt}\selectfont \centering
                \color{eptype}HISTOIRE\color{black}\\
                \freshbot\fontsize{72pt}{72pt}\selectfont
                \color{epnum}30\\[6pt]
                \cafeteria\fontsize{20pt}{20pt}\selectfont\color{epnum}(D-30)\color{black}
                \\[18pt]
                \color{black}
            }
             &
            {\fotskip\fontsize{18pt}{18pt}\selectfont \ruby{龍}{りゅう}と\ruby{天}{てん}\ruby{使}{し}の\ruby{再}{さい}\ruby{来}{らい}}
            \\[6pt]
             & {\cafeteria\fontsize{18pt}{18pt}\selectfont The Dragon and the Angel Return}                                                                                                        \\[4pt]
             & {\vnmsans\fontsize{18pt}{18pt}\selectfont Rồng và thiên thần trở lại}                                                                                                               \\[8pt]
             & {\caviar{\fontsize{13pt}{13pt}\selectfont Nhân vật: \emp{kanata} \emp{watame} \emp{kson} \emp{towa} \emp{luna} \emp{kuon2} \emp{kaiho2} \emp{ryuuhou2} \emp{suzuno2} \emp{naomi2}}}
            \\[8pt]
        \end{tabular}
    \end{adjustwidth}
\end{table}
%=========================================TRUYỆN CHÍNH========================================
\color{black}

\txtbx[2]{
    {\color{vol} \uline{Tóm tắt:}} Nối tiếp những chuyển biến tích cực sau chiến dịch Réunion và chuỗi ngày ghé thăm riêng rẽ, các cô gái Thế hệ thứ Tư lại rủ nhau trở về ký túc xá nhà Hikawa vào một ngày nghỉ đầu hè. Trong khi Kanata cùng Kson dọn dẹp căn phòng tầng bảy và trải lòng về sự tự do khi không còn phải gồng mình mang đôi cánh nặng nề, gian bếp tầng năm lại biến thành một sân khấu náo nhiệt với sự kết hợp ẩm thực của Watame, Suzuno cùng Ryuuhou. Sự xuất hiện bất ngờ của Towa và cô công chúa Luna mang đến những trận cười nghiêng ngả trước kỷ luật thép của Kaiho. Dưới ánh hoàng hôn rực rỡ nơi ban công lộng gió, sợi dây liên kết giữa những người bạn thân năm nào lại càng thêm bền chặt, chứng minh rằng dù diện mạo hay con đường có đổi thay, tình bạn của họ vẫn mãi là bến đỗ bình yên nhất.
}

\pov{Amane Kanata}

Chuyến tàu điện cao tốc tuyến Chuo hướng về phía tây lăn bánh nhịp nhàng trên những thanh ray bằng thép, bỏ lại sau lưng nhịp sống náo nhiệt, ồn ã của trung tâm Tokyo.

Tôi tựa trán vào ô cửa kính mát lạnh, đưa mắt ngắm nhìn khung cảnh bên ngoài dần chuyển dịch từ những khối bê tông chọc trời sang những rặng cây xanh mướt và những mái nhà ống đặc trưng của vùng Kawachi, thủ đô Etsunan.

Đây không còn là lần đầu tiên tôi đặt chân lên con đường này nữa.

Kể từ sau buổi chiều tái ngộ đầy nước mắt tại kho hàng ngoại ô của mẹ anh Kuon, rồi đến những chuyến ghé thăm riêng tư sau đó của từng người trong nhóm, ngôi nhà tại phường Anwa đã trở thành một điểm hẹn quen thuộc đối với cả năm đứa chúng tôi. Chúng tôi đều ghi nhớ nằm lòng những quy ước bảo mật nghiêm ngặt mà anh Kuon và chị Kaiho đã dặn dò: không bao giờ đi chung cả nhóm quá đông, không nhắc tên địa điểm trên sóng truyền hình, và luôn chú ý quan sát xung quanh mỗi khi rời ga tàu.

Hôm nay là một ngày Thứ Bảy cuối tháng Sáu nắng đẹp hiếm hoi.

Ngồi ở chiếc ghế đối diện tôi là Watame. Nàng Cừu của chúng tôi diện một chiếc váy hoa nhí thanh lịch, trên đùi đặt ngay ngắn một chiếc giỏ mây phủ khăn carô, tay thì đang chăm chú lật giở một cuốn sổ tay bìa da nhỏ nhắn ghi chép đầy những công thức nấu ăn mà bà ấy đã sưu tầm suốt tuần qua.

``Bà đang xem công thức món gì đấy, Watame?'' tôi tò mò rướn người hỏi.

Watame ngẩng đầu lên, đôi má bầu bĩnh khẽ ửng hồng cùng nụ cười tít cả mắt: ``À, tôi đang tính nhờ Suzuno-chan chỉ cho cách nấu món súp rau củ mùa hè thanh mát ấy mà! Nghe Kuon-senpai kể em ấy có kỹ thuật hầm nước dùng từ củ quả ngọt tự nhiên cực kỳ đỉnh, tôi muốn học lỏm một chút về nấu cho mấy bà thưởng thức!''

``Bà định nấu súp rau củ thật đấy à?'' tôi nhướng mày trêu chọc. ``Tôi lại cứ tưởng bà sang đấy để làm một bữa linh đình cho Kson nhắm trà chứ?''

``Này nhé! Đừng có trêu tôi thế chứ, Kanata!'' Watame phồng má hờn dỗi, hai tay ôm chặt lấy giỏ mây. ``Hôm nay tôi là đầu bếp nghiêm túc đấy nhé! Với lại dạo này Kson bận rộn nhiều việc ở xưởng cơ khí và kênh cá nhân, tôi muốn nấu món gì đó thanh đạm tẩm bổ cho bà ấy một chút thôi mà.''

Nhìn ánh mắt chân thành của Watame, tôi khẽ mỉm cười. Đúng là dù có qua bao nhiêu biến cố, tấm lòng ấm áp như len mùa đông của cô bạn này vẫn chẳng hề thay đổi chút nào.

\ct

Bước ra khỏi nhà ga Kyuukami, làn gió đầu hạ mát rượi thổi qua vạt áo làm dịu hẳn cái nóng ban trưa.

Chúng tôi rảo bước trên con dốc quen thuộc dẫn về ký túc xá. Cánh cổng gỗ lớn của nhà Hikawa mở hé như thể đã biết trước chúng tôi sẽ đến.

Vừa bước chân vào tiền sảnh tầng một, không gian yên tĩnh mát mẻ cùng mùi hương gỗ thông đặc trưng đã ùa tới chào đón. Anh Kuon đang ngồi kiểm tra sổ sách sau bàn làm việc ở tiền sảnh, trên người vẫn là chiếc áo thun đỏ giản dị quen thuộc.

``Chào hai đứa,'' anh Kuon ngẩng lên, nụ cười điềm đạm nở trên môi. ``Hôm nay đi chung tàu với nhau à? Trông tươi tỉnh thế này là biết tuần qua làm việc suôn sẻ rồi.''

``Em chào Kuon-senpai ạ!'' Watame giơ chiếc giỏ mây lên lắc lắc. ``Em có mang ít nguyên liệu tươi từ chợ nông sản sang đây này anh! Suzuno-chan có ở nhà không ạ?''

``Có, con bé đang ở gian bếp tầng năm dọn dẹp đấy,'' anh chỉ tay lên cầu thang. ``Lên đó đi, nghe bảo Wata-chan sang là Suzuno đã chuẩn bị sẵn mấy loại nấm rừng thu hoạch từ sáng sớm rồi.''

``Oa, tuyệt quá! Em xin phép lên trước nhé, Kanata!'' Watame reo lên đầy hào hứng, nhanh thoăn thoắt ôm giỏ đồ chạy tót lên cầu thang như một chú cừu nhỏ tìm thấy đồng cỏ non.

Tôi bật cười nhìn theo bóng lưng vội vã của cô bạn, rồi quay sang hỏi anh Kuon: ``Còn Kson thì sao anh? Cậu ấy đang ở trên tầng bảy ạ?''

``Ừ, đang ở trên phòng đấy,'' anh Kuon nhấp một ngụm trà, khóe mắt nheo lại vẻ hóm hỉnh. ``Sáng nay xưởng cơ khí bên kia thành phố nghỉ làm nửa ngày, nó không phải đèo xe máy sang phụ bố nên đang tranh thủ tổng vệ sinh lại phòng stream và giá sách. Em lên phụ một tay đi, con Rồng ấy mà dọn dẹp thì chẳng khác nào bão quét qua chuồng gà đâu.''

``Vâng, để em lên xem sao!'' tôi gật đầu chào anh rồi rảo bước lên tầng bảy.

\ct

Cánh cửa phòng ngoài cùng bên phải của tầng bảy khép hờ.

Tôi khẽ đẩy cửa bước vào, định cất tiếng chào thì một tiếng thở dài thườn thượt vang lên từ phía góc phòng: ``Trời đất ơi\dots\ sao mình lại nhét lắm thứ linh tinh vào cái hộc tủ này thế không biết\dots''

Kson đang ngồi bệt dưới sàn nhà gỗ, mặc chiếc áo ba lỗ màu xám bó sát cùng chiếc quần túi hộp rộng thùng thình, mái tóc tím than dài được búi lỏm xỏm trên đỉnh đầu. Xung quanh bà ấy là hàng chục thùng các-tông mở tung, ngổn ngang dây cáp nối, tay cầm chơi game cũ, cùng vô số tập tài liệu và đĩa game chất đống như núi.

Tôi khoanh tay dựa vào khung cửa, cất giọng trêu chọc: ``Này Tổng trưởng, bà đang dọn phòng hay đang chuẩn bị mở bãi đồ cũ đấy?''

Kson giật bắn mình, quay phắt đầu lại. Nhìn thấy tôi, đôi mắt tím của cô nàng lập tức bừng sáng, khóe môi nhếch lên nụ cười sảng khoái quen thuộc: ``Ủa, Kanatan đấy à?! Đến từ lúc nào mà không phát ra tiếng động gì thế, định dọa người ta đấy à?''

``Tại bà mải vùi đầu vào đống đồ cũ của bà quá chứ ai!'' tôi bật cười, bước vào phòng rồi ngồi xổm xuống cạnh một thùng các-tông. ``Để tôi phụ một tay cho, nhìn bà loay hoay một mình chắc đến tối cũng chưa xong.''

``Thế thì đội ơn thiên sứ nhé!'' Kson cười khà khà, tiện tay quẳng cho tôi một cuộn băng dính to bản. ``Phân loại giúp tôi đống đĩa game bên kia đi, cái nào của hệ máy cũ thì xếp riêng vào thùng màu xanh nhé.''

Hai chúng tôi bắt tay vào việc. Vừa lau dọn bụi bặm, chúng tôi vừa trò chuyện rôm rả về những chuyện xảy ra suốt mấy tuần qua. Kson kể về mấy pha xử lý cồng kềnh trong buổi stream game bắn súng hôm trước, còn tôi thì kể về lịch trình tập nhảy căng thẳng cho buổi hòa nhạc sắp tới.

Trong lúc dọn dẹp, ánh mắt tôi vô tình lướt qua chiếc kệ gỗ ở góc phòng.

Chiếc áo phông trắng ``PP-Tenshi'' huyền thoại vẫn nằm ngay ngắn trong khung kính acrylic, bên cạnh là dàn standee và những món đồ lưu niệm in hình tôi được lau chùi sạch bóng, không dính một hạt bụi nào. Tôi khẽ nhếch mép, liếc xéo bà ấy:

``Này, góc này vẫn được chăm sóc kỹ lưỡng quá nhỉ? Không sợ người khác bước vào phòng lại trêu à?''

Kson khựng lại một nhịp, đôi má hơi ửng đỏ, nhưng bà ấy lập tức hất cằm, vờ gắt gỏng: ``Kệ tôi chứ! Tiền của tôi, phòng của tôi, tôi thích trưng bày cái gì là quyền của tôi! Bà lắm chuyện vừa thôi nhé, coi chừng tôi bắt bà ký thêm mười chữ ký lên đống đĩa CD kia bây giờ!''

``Thôi xin, chữ ký của tôi đắt giá lắm đấy nhé!'' tôi cười lớn, né cú huých vai của cô bạn.

Dọn dẹp được chừng một tiếng đồng hồ, căn phòng cuối cùng cũng đâu vào đấy. Những chiếc thùng các-tông được xếp gọn gàng vào góc tường, mặt sàn gỗ được lau sạch bóng loáng.

Kson với tay lấy hai lon trà lúa mạch ướp lạnh từ chiếc tủ lạnh mini ở góc phòng, ném cho tôi một lon: ``Uống đi cho mát, bà tướng.''

``Cảm ơn bà nhé.''

Hai đứa chúng tôi ngồi bệt xuống tấm thảm lông cạnh ban công, bật nắp lon nước rồi nhấp một ngụm dài. Hơi lạnh ngọt mát lan tỏa khắp cổ họng, xua tan đi mọi mệt mỏi.

Tôi ngắm nhìn Kson. Dưới ánh nắng vàng rực rỡ chiếu qua khung cửa kính, những giọt mồ hôi lấp lánh trên vầng trán cô bạn. Nụ cười của bà ấy lúc này thật sảng khoái, bờ vai buông lỏng tự nhiên, không còn chút căng thẳng hay phòng thủ nào như những ngày tháng giông bão cuối năm 2021.

Một cảm giác bồi hồi bất chợt dâng lên trong lòng tôi.

``Kson này\dots'' tôi khẽ lên tiếng, giọng nói trầm xuống một nhịp.

``Hử? Sao thế?'' Kson quay sang, ngơ ngác nhìn tôi.

``Dạo này nhìn bà vui hơn hẳn đấy,'' tôi mỉm cười, nhìn thẳng vào mắt cô bạn. ``Không chỉ là vui vẻ lúc stream đâu, mà là cái thần thái này này\dots\ Nhẹ nhõm, tự do, và rạng rỡ hơn rất nhiều so với những ngày cuối cùng hồi đó.''

Kson im lặng một thoáng. Bà ấy đưa lon trà lên nhấp một ngụm nhỏ, đôi mắt nhìn ra khoảng trời xanh ngắt ngoài ban công, nơi những cánh chim chao liệng giữa nền mây trắng. Khóe môi cô bạn khẽ cong lên một nụ cười ấm áp: ``Ừ\dots\ Bà nhận ra à?''

``Bạn bè sống chung với nhau bao lâu, sao tôi không nhận ra được chứ.''

Kson đặt lon nước xuống sàn, hai tay chống ra sau lưng, ngửa đầu hít một hơi thật sâu: ``Bà biết không, Kanatan\dots\ Trước đây, tôi luôn nghĩ mình phải dang rộng đôi cánh gồng gánh che chở cho tất cả mọi người, lúc nào cũng phải thật mạnh mẽ. Càng cố, đôi vai lại càng nặng trĩu vì những quy tắc, những áp lực và kỳ vọng từ đủ phía.''

Bà ấy quay sang nhìn tôi, ánh mắt chan chứa sự thanh thản: ``Nhưng ở đây, trong căn nhà này, giữa những con người mộc mạc này\dots\ tôi được là chính mình. Tôi có thể mặc áo ba lỗ đi dép lê, có thể thoải mái cười nói, có thể phụ vác đá lát với bác Kaien, và có thể nói bất cứ điều gì mình nghĩ. Không cần phải bay cao nữa, Kanatan à. Sống nhẹ nhõm cùng mấy bà như thế này\dots\ mới thực sự là điều tôi tìm kiếm suốt bấy lâu.''

Nghe những lời chân thành ấy từ đáy lòng người bạn thân, cổ họng tôi bỗng nghẹn lại. Một luồng ấm áp lan tỏa khắp lồng ngực.

Tôi đưa tay ra, đấm nhẹ một cú vào bắp tay săn chắc của Kson: ``Tốt lắm. Nghe bà nói thế là tôi yên tâm rồi.''

Kson bật cười ha hả, quàng tay qua cổ tôi siết nhẹ: ``Này nhé, ai cho bà làm cái mặt bà cụ non đấy hả?! Coi chừng tôi kẹp cổ đo lực tay bây giờ!''

``Bà thử xem ai khỏe hơn! Lực nắm năm mươi cân của tôi không ngán bà đâu nhé!''

Hai đứa lại lao vào vật lộn trên sàn nhà như hai đứa trẻ con, tiếng cười đùa vang vọng khắp căn phòng ngập nắng.

\ct

Đến khoảng hai giờ chiều, cái bụng của cả hai bắt đầu réo lên ùng ục. Tôi và Kson quyết định đi xuống tầng năm xem tình hình ``dự án ẩm thực'' của Watame ra sao.

Thế nhưng, vừa bước chân tới chiếu nghỉ cầu thang tầng sáu, một mùi thơm nồng nặc - pha lẫn mùi thơm ngọt của súp hầm và mùi cay nồng xộc thẳng lên mũi - đã ùa ra ngoài hành lang. Đi kèm với đó là những âm thanh hỗn loạn như một bãi chiến trường.

``Ryuuhou-chan! Bỏ lọ ớt bột đó xuống ngay! Nồi súp này là vị thanh ngọt dịu mát mà, ai lại cho cả đống ớt vào thế kia!'' tiếng Watame kêu lên đầy hoảng hốt.

``Thanh ngọt ăn mãi chán lắm Watame-san! Trời nóng thế này là phải ăn cay một chút, toát mồ hôi ra mới đã chứ! Thêm cả sốt mayonnaise nữa cho nó béo ngậy!'' giọng cô em gái của cậu Tổng cãi lại đầy bướng bỉnh.

``Hai người bình tĩnh đã ạ\dots\ Nếu cho thêm mayonnaise lúc này thì chất béo sẽ bị tách lớp, nồi súp sẽ bị váng dầu đấy ạ\dots'' giọng Suzuno đều đều, cố gắng dùng lý lẽ hóa học để can ngăn hai người.

Tôi và Kson bước vào bếp, nhìn thấy cảnh tượng trước mắt mà không nhịn được cười.

Watame đang một tay giữ chặt nắp nồi súp, tay kia ra sức ngăn cản Ryuuhou đang cầm lọ ớt bột đỏ rực lăm lăm xông tới. Suzuno thì đứng bên cạnh, hai tay cầm hai chiếc đũa dài với vẻ mặt vô cùng bối rối. Cả gian bếp lúc này chẳng khác nào một bãi chiến trường ẩm thực nho nhỏ.

``Mấy đứa đang nấu súp hay đang làm thí nghiệm phá hoại thế hả?'' Kson bước tới, giơ tay giật phắt lọ ớt bột từ tay Ryuuhou, cốc nhẹ lên đầu cô bé một cái. ``Nhóc con, muốn phá bếp để chị Kaiho lên phạt dọn sân à?''

Ryuuhou xoa đầu, bĩu môi lầm bầm: ``Kson-san lúc nào cũng bênh khách\dots''

Watame thở phào một hơi nhẹ nhõm, quệt mồ hôi trên trán rồi cười tươi rói khi thấy tôi: ``May quá, hai bà xuống rồi! Súp hầm rau củ của tôi sắp xong rồi đấy, chỉ cần đun nhỏ lửa thêm mười phút nữa cho ngấm thôi!''

Nói rồi, cô bạn tháo tạp dề, nhẹ nhàng khuấy chiếc muôi gỗ trong nồi. Mùi hương ngọt lịm của bắp non, cà rốt và nấm rừng lan tỏa khắp không gian, khiến cơn đói của chúng tôi càng thêm cồn cào.

Trong lúc chờ đợi súp sôi liu riu, Watame khẽ đung đưa người theo một nhịp điệu vô hình. Rồi bất giác, cô ấy cất tiếng hát vu vơ. Đó là một khúc hát êm ả, mộc mạc và ấm áp quen thuộc mà cô ấy vẫn thường ngâm nga mỗi khi phát sóng trực tiếp. Giọng hát của Watame tựa như một tấm chăn len mềm mại, xoa dịu đi mọi sự ồn ào vừa rồi, mang lại một cảm giác bình yên đến lạ lùng cho cả gian phòng.

Ryuuhou ngừng càu nhàu, ngơ ngẩn đứng tựa vào bồn rửa lắng nghe.

Từ phòng khách phía ngoài, hai bóng dáng nhỏ nhắn rón rén bước vào.

Là bé Naomi đang ôm chú gấu bông tròn ủng, cùng với bạn Koneko đeo chiếc túi nhỏ sau lưng.

Naomi chớp chớp đôi mắt to tròn trong veo, ngước nhìn Watame như nhìn thấy một cô tiên bước ra từ trang sách cổ tích. Naomi kéo chiếc ghế gỗ nhỏ ngồi xuống góc bàn cùng Koneko, hai đứa nhỏ lặng lẽ lắng nghe từng nốt nhạc trong trẻo.

Kuon từ cửa bước vào, thấy cảnh tượng ấy bèn khẽ ra hiệu cho chúng tôi giữ im lặng. Anh đứng tựa vào khung cửa, mỉm cười ấm áp nhìn hai cô bé đang say sưa đắm chìm trong giai điệu của người bạn cũ.

Khoảnh khắc ấy, gian bếp không chỉ ngập tràn mùi hương của món súp ngon lành, mà còn đong đầy hơi ấm của sự gắn kết gia đình.

\ct

Khoảng ba giờ chiều, khi nồi súp vừa được tắt bếp thì ở bên ngoài hành lang, tiếng bước chân huỳnh huỵch cùng tiếng cãi cọ quen thuộc đã vang lên tận hành lang tầng năm.

``Đã bảo là đi thang máy rồi mà cứ bắt leo thang bộ hoài à! Mỏi hết cả chân rùi nè\dots''

``Tại bà cứ nằng nặc đòi ghé tiệm kem làm trễ chuyến tàu đấy chứ! Mau vào đi, đứng càu nhàu ở cửa làm gì!''

Cánh cửa phòng khách bật mở.

Towa và Luna xuất hiện!

Towa diện một bộ đồ thể thao đen phối tím năng động, một bên vai đeo chiếc túi vải đựng tài liệu tập luyện. Còn Luna mặc chiếc váy xòe màu hồng phấn dễ thương, hai má phồng to vẻ hờn dỗi. Thì ra sau khi kết thúc buổi tập duyệt riêng ở công ty, hai người cũng rủ nhau bắt chuyến tàu tiếp theo sang đây tụ họp.

``A! Mọi người ở đây hết rồi này!'' Towa tròn mắt ngạc nhiên khi thấy đông đủ cả bọn.

``Towa-chan! Luna-chan!'' Watame mừng rỡ vẫy tay. ``Hai bà tới đúng lúc lắm, súp vừa nấu xong luôn này!''

Vừa nhìn thấy gian phòng khách ấm cúng, Luna lập tức thả mình ngồi phịch xuống chiếc ghế sofa êm ái nhất, vươn vai một cái thật dài: ``Oa\~{} Mệt xỉu luôn á! Có trà lạnh hay bánh ngọt gì giải khát không mọi người ơi\~{}?''

Ryuuhou đang bưng khay bát đĩa ra bàn, nghe thấy thế bèn liếc nhìn rồi trêu chọc: ``Ở cái nhà này không có người hầu đâu nhé Luna-san! Muốn ăn thì lát nữa lăn vào bếp mà bưng bát đĩa ra nha!''

Luna lập tức phụng phịu, chỉ tay về phía Kuon: ``Kuon-senpai ơi, coi em gái anh bắt nạt em kìa!''

Điệu bộ giả vờ mếu máo của Luna làm cả bọn chúng tôi cười lăn cười bò.

Thế nhưng, niềm vui chưa kéo dài được bao lâu thì tiếng bước chân nặng trịch cùng tiếng kim loại va chạm lách cách bất chợt vang lên từ cầu thang dẫn lên tầng năm.

Hóa ra là chị Kaiho vừa đi xe máy từ xưởng cơ khí cách bốn phường về nhà sớm hơn thường lệ. Vừa dựng xe ở sân trước bước vào nhà, tiếng ồn ào nhõng nhẽo của Luna trên gác đã vọng xuống tận tiền sảnh, khiến chị tiện tay cầm luôn chiếc cờ lê kim loại dài nửa mét vừa bảo dưỡng máy xong bước thẳng lên kiểm tra.

*RẦM!*

Cánh cửa phòng khách bị đẩy mạnh sang một bên. Chị Kaiho đứng sững ở ngưỡng cửa.

Trên người chị vẫn là bộ đồ yếm bảo hộ lao động màu xám lấm lem vài vệt dầu máy, chiếc cờ lê kim loại khổng lồ cầm hờ hững trên tay phải. Đôi mắt xanh biếc sắc bén quét một vòng qua căn phòng khách đang nhốn nháo.

Luna lập tức nín bặt, đôi mắt mở to nhìn người phụ nữ to lớn đầy uy quyền trước mặt.

Kaiho bước tới hai bước, chiếc cờ lê gõ nhẹ một tiếng *CẠCH* dứt khoát lên mặt bàn gỗ: ``Đủ rồi đấy nhé.''

Giọng nói của chị trầm thấp nhưng mang sức nặng: ``Ở cái ký túc xá này, tất cả mọi người đều như nhau. Dù là ai đi chăng nữa, hễ bước chân qua cánh cửa kia thì đều là thành viên trong nhà. Ăn xong phải tự mang bát đĩa đi rửa, rác phải tự phân loại, không có ai ngoại lệ cả. Hiểu rõ chưa?''

Toàn bộ phòng khách rơi vào sự im lặng.

Towa nuốt nước bọt đánh ực, vội vàng rụt cổ lại. Kson và Watame thì mím môi nhịn cười.

Còn Luna\dots\ cô nàng ngước nhìn bờ vai rộng vững chãi của Kaiho, nhìn chiếc cờ lê khổng lồ trong tay chị, đôi mắt to tròn bỗng chớp chớp vẻ bái phục tột cùng. Sự nhõng nhẽo ban nãy biến mất, thay vào đó là giọng nói lí nhí, ngoan ngoãn đến lạ kỳ: ``D\dots\ dạ vâng ạ\dots\ Chị\dots\ chị ngầu quá đi mất\dots''

Cả căn phòng bỗng chốc vỡ òa trong những tràng cười không dứt. Ngay cả Kaiho sau khi nghe câu khen ngợi ngây ngô ấy cũng không kìm được mà khẽ nhếch khóe môi, quay người bước về phía bồn rửa tay: ``Biết thế là tốt. Mau dọn bát đĩa ra ăn đi, súp nguội hết bây giờ.''

\ct

Bữa ăn trưa muộn diễn ra trong sự ấm cúng và rộn rã hiếm có.

Món súp rau củ của Watame và Suzuno nhận được cơn mưa lời khen từ tất cả mọi người. Vị ngọt thanh của ngô non hòa quyện với vị bùi bùi của nấm rừng tạo nên một hương vị thanh khiết, làm dịu đi cái nóng oi ả của buổi trưa hè.

Cả năm đứa chúng tôi - nhóm \textit{holoForce} và Kson - quây quần bên chiếc bàn gỗ lớn, vừa ăn vừa tranh nhau kể chuyện, thi thoảng lại cãi cọ chí choé như những ngày tháng còn ở chung dưới một mái nhà.

Đến khoảng năm giờ chiều, khi ráng chiều bắt đầu nhuốm màu đỏ ối lên nền trời phía tây, tôi và Kson lại rủ nhau bước ra ngoài ban công tầng bảy.

Gió chiều từ những rặng núi xa xôi thổi về lồng lộng, mang theo hơi thở mát lành của cỏ cây.

Tôi tựa hai cánh tay vào lan can kim loại, phóng tầm mắt nhìn về phía thị trấn Kawachi đang dần chìm trong ánh hoàng hôn tuyệt đẹp. Từng dải mây màu cam rực rỡ đan xen với sắc tím thẫm tạo nên một bức tranh thiên nhiên kỳ ảo.

Kson đứng cạnh tôi, hai tay đút túi quần, mái tóc tím than khẽ bay trong gió.

``Tuyệt thật đấy nhỉ,'' Kson khẽ thở dài một tiếng đầy thỏa mãn.

Tôi quay sang nhìn cô bạn thân, ngắm nhìn sườn mặt góc cạnh đầy bản lĩnh của bà ấy dưới ánh ráng chiều. Một cảm xúc sâu lắng, nghẹn ngào bỗng dâng lên nơi cuống họng.

Tôi khẽ thì thầm, gọi lại cái tên thân thương đã khắc sâu trong tim suốt bao năm tháng: ``Coco này\dots''

Kson thoáng khựng lại một nhịp. Bà ấy không quay đầu sang, nhưng tôi biết bà ấy đang lắng nghe từng lời của tôi.

``Tôi thật sự rất mừng\dots\ vì bà vẫn ổn.''

Giọng tôi run run, chứa chan tất cả sự chân thành và những lo âu bị dồn nén suốt một năm qua: ``Nhìn thấy bà mỉm cười như hôm nay, nhìn thấy bà tìm được một bến đỗ bình yên như thế này\dots\ Tôi mới thực sự trút bỏ được gánh nặng trong lòng. Bọn tôi nhớ bà nhiều lắm, nhưng nhìn thấy bà hạnh phúc thế này\dots\ bọn tôi an tâm rồi.''

Một khoảng lặng êm đềm trôi qua giữa hai chúng tôi. Chỉ có tiếng gió rì rào qua tán lá cây phong dưới sân và tiếng cười đùa ríu rít của bọn Watame vọng lên từ tầng dưới.

Kson từ từ quay sang nhìn tôi. Trong đôi mắt tím của cô bạn, tôi nhìn thấy những tia sáng lấp lánh như ngàn vì sao đêm.

Bà ấy giơ một tay lên, xoa mạnh mái tóc ngắn của tôi làm nó rối tung lên, khóe miệng cong lên một nụ cười rạng rỡ và ấm áp nhất trần đời: ``Ừ\dots\ Ổn hơn rất nhiều rồi, Kanatan ạ.''

Kson nhìn xuống khoảng sân ký túc xá ngập tràn ánh nắng chiều, giọng nói trầm ấm lan tỏa trong không gian:

``Nhờ có cái gia đình kỳ lạ này che chở\dots\ và quan trọng nhất là nhờ có mấy bà luôn dõi theo tôi đấy.''

Tôi bật cười, giơ nắm đấm huých nhẹ vào sườn cô bạn một cái.

Hai đứa chúng tôi cùng tựa vào lan can, ngắm nhìn mặt trời từ từ lặn xuống sau rặng núi xa xôi.

Dù thời gian có trôi đi, dù diện mạo hay tên gọi có đổi thay, và dù mỗi người chúng tôi đang bước đi trên những ngả đường riêng biệt\dots\ sợi dây gắn kết của Thế hệ thứ Tư sẽ mãi mãi không bao giờ đứt gãy.

Bởi vì chúng tôi đã luôn là gia đình của nhau.

%==========================================TRUYỆN PHỤ=========================================
\omk

\pov{-}

Khoảng bốn giờ chiều hôm đó, dưới phòng khách tầng năm.

Mọi người đang bận rộn thu dọn bát đĩa sau bữa ăn, chỉ có Towa là đang ngồi ở góc bàn với vẻ mặt vô cùng đăm chiêu.

Cô nàng khoanh hai tay trước ngực, cố giữ vẻ mặt nghiêm trang: ``Mọi người đừng có hiểu lầm nhé! Hôm nay tôi chỉ tiện đường sau buổi tập thì ghé qua kiểm tra tình hình thôi đấy. Ai ngờ mấy bà kéo đến đông đủ thế này.''

Watame vừa lau bàn vừa khúc khích cười: ``Vâng vâng, Towa-chan bận rộn ghé qua thăm là tụi này mừng lắm rồi ạ\~{}''

``Này! Tôi nói thật mà!'' Towa đỏ mặt phân bua.

Đúng lúc đó, Kuon từ phòng làm việc bước ra, trên tay cầm chiếc kìm nhỏ và cuộn chỉ dù: ``Towa-chan, cái quai túi xách của em bị đứt chỉ hồi nãy anh khâu lại xong rồi này. Thử đeo lại xem đã vừa tầm chưa?''

``A\dots\ dạ! Em cảm ơn Kuon-senpai ạ!''

Towa lập tức đứng dậy, lon ton chạy lại chỗ Kuon. Cô nàng nhận lấy chiếc túi, kiểm tra đường chỉ khâu tỉ mỉ và chắc chắn của anh Tổng quản lý, đôi mắt bừng sáng vẻ thán phục: ``Oa, anh khâu khéo tay thật đấy ạ! Chẳng bù cho em, cầm kim chỉ toàn tự đâm vào tay thôi\dots''

Kuon mỉm cười, ngồi xổm xuống cạnh chiếc ghế để chỉnh lại độ dài của dây đeo cho vừa vặn với vóc dáng của cô bé: ``Không có gì đâu, hồi ở ký túc xá cũ anh phải tự sửa đồ cho mấy đứa suốt ấy mà. Đứng yên anh chỉnh lại nấc cài cho nhé.''

Towa ngoan ngoãn đứng yên. Nhìn góc nghiêng nghiêm túc và ân cần của người quản lý đã luôn đồng hành cùng mình từ những ngày đầu ra mắt, khóe môi cô bé bất giác cong lên một nụ cười ngượng ngùng.

Towa lén thò tay vào túi áo khoác, nhẹ nhàng rút chiếc điện thoại thông minh ra.

Cô nàng giơ điện thoại lên ngang ngực, nghiêng nhẹ góc máy sao cho bắt trọn được gương mặt rạng rỡ của mình và góc nghiêng đang chăm chú sửa túi của Kuon ở phía sau, ngón tay chuẩn bị bấm nút chụp màn hình\dots

*CẠCH*

``Úi chà chà! Bắt quả tang tại trận nhé!'' Một giọng nói lanh lảnh bất ngờ vang lên ngay sát vành tai Towa.

Ryuuhou từ sau cánh cửa bếp nhảy bổ ra, hai mắt sáng rực như phát hiện ra kho báu, chỉ thẳng tay vào màn hình điện thoại của cô nàng: ``Towa-san! Chị vừa lén bật camera trước chụp ảnh selfie với Onii nhà em đúng không?!''

Towa giật thót tim như bị điện giật, chiếc điện thoại suýt chút nữa thì rơi thẳng xuống đất. Toàn bộ khuôn mặt cô bé từ trán xuống tận mang tai đỏ bừng lên như một trái cà chua chín mọng. Cô vội vàng giấu phắt chiếc điện thoại ra sau lưng, cuống cuồng lắc đầu xua tay lia lịa: ``\textbf{K-Không có! Đừng có mà vu oan!}''

``Có mà rành rành ra kia kìa! Em thấy rõ ràng góc máy có cả mặt anh ấy luôn nhé!'' Ryuuhou khoanh tay cười khoái chí, quay sang réo gọi cả nhà: ``Mọi người ơi ra mà xem nè! Towa vừa lén chụp ảnh kỷ niệm với anh trai em này!''

Cả phòng khách lập tức ồ lên thích thú. Kson huýt sáo một tiếng giòn giã, còn Watame và Luna thì ôm bụng cười ngặt nghẽo.

Towa tức tối dậm chân thình thịch xuống sàn, giọng nói lạc hẳn đi vì xấu hổ: ``Em\dots\ em chỉ đang đo góc sáng để thử camera thôi mà! Đâu có lén lút chụp ảnh gì đâu chứ!''

Kuon ngơ ngác ngẩng đầu lên, tay vẫn đang cầm chiếc quai túi xách, gãi gãi đầu đầy vẻ khó hiểu:

``Hử? Chụp ảnh gì cơ? Có chuyện gì thế mấy đứa?''

Nhìn vẻ mặt ngơ ngác của anh Tổng quản lý và điệu bộ cuống cuồng muốn độn thổ của Towa, cả gian phòng khách nhà Hikawa lại một lần nữa bùng nổ trong những tràng cười giòn giã nhất trần đời.

\end{document}